\documentclass[a4paper,12pt]{article}
\setcounter{secnumdepth}{5}
\setcounter{tocdepth}{3}
\input{/usr/share/LaTeX-ToolKit/template.tex}
\begin{document}
\title{元素週期性}
\author{沈威宇}
\date{\temtoday}
\titletocdoc
\section{元素週期性（Periodic Trends）}
元素週期律/（元素）週期性（趨勢）（Periodic Trends）
\subsection{簡史}
\sssc{德貝萊納三元素組（Döbereiner's triads）}
1829年德貝萊納提出，列出四組每組三個化性相似的元素：(Ca, Sr, Ba)、(Cl, Br, I)、(S, Se, Te)、(Li, Na, K)。
\subsubsection{紐蘭的八度律（Law of Octaves）}
1864年紐蘭提出，認為元素的原子量每相差八，其化學性質相似。
\subsubsection{門得列夫的原子量元素週期表}
1869年門得列夫提出，認為元素的性質隨原子量的增大週期性重複，並依照原子量排列當時已知的六十三種元素，並成功預言部分元素，但氬、鉀、鈷、鎳、碲、碘的原子量所示位置與化學性質不符。
\subsubsection{莫斯利的 X 射線光譜頻率與原子序正比與原子序元素週期表}
1914年，莫斯利利用 X 射線晶體繞射法觀察和測量了多種金屬化學元素的光譜，發現莫斯利定律，故認為原子序才是決定元素化學性質的主要因素，並以之排列元素週期表。 
\ssc{現在的元素週期表（Periodic Table）}
根據原子序排列，縱行稱族（Group）、橫列稱週期（Period），共7週期、18族、118元素，如下：
\bctf\icg[width=1.0\tw]{Cambridge.jpg}\cpt{Cambridge University Press. 2014. https://doi.org/10.1017/CBO9781139096782.017.}\efct
\bctf\icg[width=1.0\tw]{956-2018_orig.jpg}\cpt{LiFe 生活化學. 2018. Emperor Chemical Co., Ltd. https://www.lifechem.tw/periodictable.html.}\efct
\sssc{週期（Period）}
週期數即最外主殼層主量子數。
\sssc{族（Group）}
\bit
\item 新式分類法：IUPAC 制定，由左至右編號1至18。
\item 舊式分類法：新式分類法的1到18族依序為 IA, IIA, IIIB, IVB, VB, VIB, VIIB, VIIIB, VIIIB, VIIIB, IB, IIB, IIIA, IVA, VA, VIA, VIIA, VIIIA 族。
\item 內過渡元素不適用族。
\eit
\sssc{分區（Block）}
令週期數$n$。
\begin{itemize}
\item\tb{s 區（s-block）}：指遞建原理中最後填充電子的副殼層為 $n$s 者，即第一、二族及氦。
\item\tb{p 區（p-block）}：指遞建原理中最後填充電子的副殼層為 $n$p 者，即第十三到十八族除了氦。
\item\tb{d 區（d-block）}：指遞建原理中最後填充電子的副殼層為 $(n-1)$d 者。
\item\tb{f 區（f-block）}：指遞建原理中最後填充電子的副殼層為 $(n-2)$f 者。
\item\tb{主族元素（Main-group elements）/典型元素（Representative elements）/A 族元素}：指 s 區與 p 區。
\item\tb{鹼金屬（Alkali metals）}：第1族除了氫。
\item\tb{鹼土金屬（Alkaline earth metals）}：第2族。
\item\tb{硼族元素（Boron group）}：第13族。
\item\tb{碳族元素（Carbon group）}：第14族。
\item\tb{氮族元素（Pnictogens）}：第15族。
\item\tb{氧族元素（Chalcogens）}：第16族。
\item\tb{鹵素（Halogens）}：第17族。
\item\tb{貴族氣體（Noble gases）/惰性氣體（Inert gases）/稀有氣體/鈍氣}：第18族。
\item\tb{過渡元素（Transition elements）/過渡金屬（Transition metals, TM）/B 族元素}：指主族元素以外的元素。門得列夫首提出此詞時代表第 8、9、10 族元素。
\item\tb{鑭系元素（Lanthanide or lanthanoid）}：指原子序 57 至 71 的元素。
\item\tb{錒系元素（Actinide or actinoid）}：指原子序 89 至 103 的元素。
\item\tb{內過渡元素（Inner transition elements）/內過渡金屬（Inner transition metals, ITM）}：指鑭系元素與錒系元素。
\item\tb{貧金屬（Poor metals）/後過渡金屬（Post-transition metals）/後過渡元素（Post-transition elements）}：指第十三族至第十七族的金屬元素。
\end{itemize}
\sssc{元素週期律/（元素）週期性（趨勢）（Periodic Trends）}
指元素的各種性質在元素週期表中存在的特定模式。

下關於各性質之元素週期性趨勢之描述多為近似，且多主要針對前五週期而言。
\ssc{價（US: Valence, UK: Valency）}
\sssc{價/價數（US: Valence, UK: Valency）/價電子（Valence electron）數}
\begin{itemize}
\item IUPAC 定義：某元素的原子或原子團能夠結合的最大一價原子（如氫或氯原子）數，或能替代該元素的一價原子數。
\item 另一個定義：某元素原子在其氫化物中結合的氫原子數，或某元素原子在其氧化物中結合的氧原子數的二倍。
\end{itemize}
\sssc{價電子（Valence electron）}
在前述價之定義中參與化學鍵的電子。
\sssc{價軌域（Valence orbital）/價副殼層}
指價電子所在的軌域種類/副殼層。
\sssc{價殼層（Valence shell）}
指價電子所在的殼層。
\sssc{元素週期性}
\bit
\item 同族：化學性質相似，除氦外價電子數相同。
\item 同週期：價殼層相同，週期數等於其最外殼層主量子數。
\item 主族元素：價電子為最外主殼層之所有電子，價電子數等於其A族數。
\item 過渡元素：價電子數無簡單方法判斷。
\end{itemize}
\subsection{原子半徑（Atomic radius）}
\bctf\icg[width=0.9\tw]{AN.png}\cpt{Dot145, 2010.}\efct
\subsubsection{原子半徑（Atomic radius）}
通常指共價半徑（covalent radius）與金屬半徑（metallic radius），分別指共價鍵或金屬鍵鍵結的兩該元素原子之原子核距離（核間距）之一半。

共價單鍵之鍵長約等於兩原子共價半徑之和。

原子半徑約 30 pm 到 3000 pm；氦原子最小，半徑約 31 pm；氫原子次之，半徑約 53 pm；元素週期表上元素以銫原子最大，半徑約 298 pm。
\sssc{離子半徑（Ionic radius）}
指離子晶體結構中單原子離子的半徑使得兩種單原子離子形成的離子晶體其相鄰兩種離子之核間距約為兩種離子之離子半徑之和。
\sssc{凡得瓦半徑（Van der Waals radius）}
元素分子晶體中相鄰未鍵結原子核間距之一半。
\sssc{測量}
原子核間距離可用X射線繞射法測得。
\subsubsection{元素週期性}
\begin{itemize}
\item 同族：隨原子序增而增，因價軌域層數增加。
\item 同週期：隨原子序增而減，因核電荷引力增加大於電子間斥力增加。過渡元素多例外，一般第十族反彈向上，第十二或十三族反彈向下。同族週期增減一的影響約為同週期減增0.5-1.5個副殼層的影響。
\item 同元素：隨核外電子數增而增，增減一個電子的影響約為同週期增減1.25-2.25個副殼層的影響。
\item 同電子組態：隨原子序增而減。
\item 略可近似為：先看週期，週期愈大半徑愈大；不同元素者次看原子序，原子序愈大半徑愈小；同元素者次看核外電子數，核外電子愈多半徑愈大。
\end{itemize}
\sssc{氟斥效應}
指第二週期半徑較小、電子間斥力較大所造成的影響。如：
\bit
\item 電子親和力：Cl>F>Br>I、S>O
\item 鍵能：\ce{Cl2}>\ce{Br2}>\ce{F2}>\ce{I2}
\eit
\sssc{原子容}
一莫耳原子在固態的體積。
\bit
\item 第一週期：\ce{He}>\ce{H2}。
\item 第二週期：8A最大、1A次之、3A最小。
\item 第三至六週期：1A最大、8A次之、3A最小。
\eit
\subsection{熔沸點}
一般為共價網狀固體>金屬、離子晶體>分子晶體，惟金屬跨度較大。

第三週期元素熔沸點：
\begin{itemize}
\item 熔點：Si>Al>Mg>Na>\ce{S8}>\ce{P4}>\ce{Cl2}>Ar
\item 沸點：Si>Al>Mg>\ce{S8}>Na>\ce{P4}>\ce{Cl2}>Ar
\end{itemize}
其中 Si 為共價網狀固體，Al、Mg、Na 為金屬固體，\ce{S8}、\ce{P4} 為分子固體，\ce{Cl2} 為雙原子分子氣體，Ar 為單原子氣體。
\subsection{游離能（Ionization energy, IE）}
\bctf\icg[width=0.9\tw]{IE.png}\caption{First ionization energies. Double sharp, 2021.}\efct
\sssc{定義}
指氣體原子或陽離子移去束縛力最弱的電子到無窮遠處所需之能量，即價軌域與無限遠處的能量差。游離前非陰離子者必大於零。
\sssc{連續游離能}
基態氣體原子，由外向內依序游離，第$i$個游離的電子的游離能為第$i$游離能$IE_i$。未聲明時，通常指$IE_1$。
\sssc{測定方法}
在陰極射線管裝入約 \scinote{6}{-2}至\scinote{2.6}{-5} atm 的低壓氣體，加以電壓，測電流。若提高電壓到某一值時電流強度劇增，該電壓稱游離電位，所相當之能量即游離能，第$i$次電流劇增為$IE_i$。
\subsubsection{影響因素}
\begin{itemize}
\item 電子離核愈遠，$IE$愈小。
\item 核電荷愈大，對電子束縛力愈大，$IE$愈大。
\item 價軌域全滿、半滿較穩定，增加$IE$。 
\item 遮蔽效應：內層電子會減弱原子對外層電子的吸引力，使$IE$降低。
\end{itemize}
\subsubsection{元素週期性}
\begin{itemize}
\item 同族：隨原子序增而減，因外層電子離核較遠，且內層電子數較多（遮蔽效應）。
\item 同週期：隨原子序增而鋸齒狀遞增，因外層電子距核距差不大，而質子數增故核電荷吸引力增。但過渡元素增勢較緩且多例外。
\item 同週期逆勢下降：同週期同淨電荷數之原子中，$i+1$個價電子者的價軌域能量高於$i$個價電子者的價軌域能量時（如因全滿、半滿），前者的游離能小於後者，即逆勢下降，如：
\begin{itemize}
\item 前三週期：二個價電子者>三個價電子者，因為前者 $n$s 全滿，較穩定。
\item 第二週期$IE_1$、$IE_2$及第三週期$IE_1$：五個價電子者>六個價電子者，因為前者 $n$p 半滿，較穩定。
\end{itemize}
逆勢下降者仍較上一週期上一族者大，如 $IE_1$：B>Mg, S>N。
\item $IE_{i+1}$對原子序之圖形略同$IE_i$對原子序之圖形向原子序正向移一單位、能量正向移一些之圖形。
\item He 為$IE_1$最大之元素，2372 kJ/mol；Cs 為自然界$IE_1$最小之元素，376 kJ/mol。
\item 第一列過渡元素：
\begin{itemize}
\item $IE_1$：Zn最大、Sc最小，Zn 遠大於 Sc 至 Cu。
\item $IE_2$：Cu最大、Sc最小。
\item $IE_3$：Zn最大、Sc最小。
\end{itemize}
\end{itemize}
\subsubsection{同一元素比較}
\begin{itemize}
\item $\forall i < j\colon IE_i < IE_j$。
\item 若 $IE_i$ 與 $IE_{i+1}$ 所游離的電子在不同副殼層，或 $IE_{i+1}$ 游離前價軌域恰全滿或半滿，則 $IE_i\ll IE_{i+1}$。
\item 因價電子數多斥力較大，若最後填入的數顆電子在同一副殼層，則游離該等電子的$IE$大約正比於游離前的$\frac{Q^+Q^-}{r}$，其中$Q^+$為質子數、$Q^-$為電子數、$r$為原子半徑。如：IIA之$\frac{IE_2}{IE_1}\approx 2$，IIIA之$\frac{IE_3}{IE_2}\approx\frac{3}{2}$，但IIIA之$\frac{IE_2}{IE_1}\gg 2$。
\end{itemize}
\subsection{電子親和力（Electoron affinity, EA）}
\bctf\icg[width=0.9\tw]{EA.png}\cpt{DePiep, 2012.}\efct
\sssc{定義}
指氣體陰離子移去束縛力最弱的電子到無窮遠處所需之能量，即價軌域與無限遠處的能量差。

電子親和力過去的定義為現今定義（以上所述之定義）之相反數，此不用。
\sssc{連續電子親和力}
一原子的第$i$電子親和力$EA_i$定義為該原子的$(-i)$價基態離子游離一個電子變成$(-i+1)$價基態離子所須吸收的能量，即第$(-i+1)$游離能$IE_{-i+1}$。未聲明時，通常指$EA_1$。
\subsubsection{影響因素、元素週期性與同一元素比較}
\begin{itemize}
\item 電子親和力$EA_i$視為$IE_{-i+1}$，影響因素與同一元素比較同$IE$，元素週期性略同$IE$但同族隨週期增電子親和力減勢較不明顯。
\item 多為正，因放出電子須吸熱，但鈹、鎂、氮、8A 為負，因價軌域全滿或半滿。
\item 無論元素，所有$EA$均小於所有$IE$；Cl 為$EA_1$最大之元素，349 kJ/mol，小於自然界$IE_1$最小之元素 Cs 之$IE_1$。
\item 第一週期$EA_1$：H>0>He
\item 第二週期$EA_1$：7A>6A>4A>1A>3A>0>5A>2A>8A
\item 第三週期$EA_1$：7A>6A>4A>5A>1A>3A>0>2A>8A
\item 第四至五週期主族元素$EA_1$：7A>6A>4A>5A>1A>3A>2A>0>8A
\item 氟斥效應，Cl>F>Br>I、S>O。
\end{itemize}
\subsection{電負度/陰電性（Electronegativity, EN）}
\bctf\icg[width=0.9\tw]{EN.jpg}\cpt{DMacks, 2013.}\efct
\sssc{鮑林電負度（Pauling electronegativity）}
1932年，鮑林依原子在多元素化合物中對共價鍵中電子之相對吸引能力提出電負度（electronegativity），設定氟的電負度為 4.0（電負度最大，最活潑非金屬），以之為基準得其他之電負度，一般不討論惰性氣體。此採之。
\sssc{馬利肯電負度（Mulliken electronegativity）}
1934年馬利肯提出，等於第一游離能與第一電子親和力之和除以二再除以一電子伏特。此不用。
\subsubsection{元素週期性}
\begin{itemize}
\item 同週期：電負度略隨原子序增而增。過渡元素較不守此原則而多相近且變化較不規則。
\item 同族：電負度略隨原子序增而減。
\item 主族元素：$n$週期$m$族元素電負度約在$(n-1)$週期$(m-2)$族與$(m-1)$族電負度之間。
\item Cs 為自然界電負度最小者，0.7，即最活潑金屬。
\item 氫電負度2.2；第二週期除氖外第$i$A族電負度約為$0.5(i+1)$。
\end{itemize}
\subsection{金屬性與非金屬性}
\sssc{定義}
\bit
\item 金屬性一般指還原性、氫化物/氧化物/氫氧化物溶於水鹼性、延展性、（因自由電子流動）、導熱性（因自由電子碰撞）等性質。
\item 非金屬性一般指氧化性、氫化物溶於水酸性、延展性差、導電性差、導熱性差等性質。
\item 該等性質有相似的元素週期性。
\eit
\sssc{元素週期性}
電負度愈大，金屬性愈弱，非金屬性愈強；電負度愈小，金屬性愈強，非金屬性愈弱。
\sssc{導電性}
金屬與石墨為導體，具高導電性。金屬導電性隨溫度上升而下降（因原子振動加劇，使自由電子不易通過），並一般隨微量雜質占比增加而下降。

NTP 導電性：Ag>Cu>Au>Al>Rh>Zn>Ni>Ca>Fe>Pt>Pd>Sn>Cr>Pb>Ti>Hg。
\sssc{類金屬（Metalloid）}
已確定為類金屬者包括硼、矽、鍺、砷、銻、碲。導電性介於金屬與非金屬間。純質導電性一般不太高，但 p 型或 n 型摻雜後導電度上升。
\sssc{酸鹼性}
\bit
\item 一元素的可溶氫氧化物或氫化物，電負度大於氫者溶液為酸，電負度小於氫者溶液為鹼，電負度愈大，溶液愈酸。
\item 一元素的可溶含氧酸，通常中心原子氧化態愈高，因中心原子電負度愈大，故溶液愈酸，但磷例外，酸性：\ce{H3PO3}>\ce{H3PO2}>\ce{H3PO4}。
\eit
\ssc{對角線關係/對角線規則（Diagonal relationship）}
指部分元素與下一週期下一族的元素的一些性質相似，具有此種關係者如：
\bit
\item 鋰與鎂：
\bit
\item 磷酸鋰與磷酸鎂難溶於水，碳酸鋰、碳酸鎂微溶於水，其它鹼金屬碳酸鹽和磷酸鹽則易溶於水。
\item 碳酸鋰與碳酸鎂受熱分解成氧化物與二氧化碳，其它鹼金屬碳酸鹽則需極高溫度才分解。
\item 鋰與鹼土金屬可在氮氣中燃燒生成氮化物，其它鹼金屬則否。
\item 鋰與鎂可以形成較穩定的共價有機金屬化合物，其他鹼金屬與鹼土金屬則將極不穩定。
\item 氯化鋰與氯化鎂具有強吸水性，其他鹼金屬氯化物則弱。
\eit
\item 鈹與鋁：許多鈹與鋁的化合物具有共價性。
\item 硼與矽：
\bit
\item 元素態為類金屬，其他硼族元素則否。
\item 三氧化二硼與二氧化矽均為共價鍵結、無定形態均穩定，其他硼族元素則否。
\item 硼烷和矽烷穩定者多，其他硼族元素則否。
\eit
\item 碳與磷：
\bit
\item 元素態有多種同素異形體。
\item 甲酸酸性強於碳酸，且具還原性可被氧化成二氧化碳；亞磷酸與次磷酸酸性強於磷酸，且具還原性分別可被氧化成磷酸與亞磷酸。
\eit
\eit
\end{document}
